\chapter{Nonadic Return, Memory, and Cartan Holonomy}
\label{chap:puente-retorno-holonomia-cartan}
\index{return!27--54--108}\index{Holonomy!TPK--Cartan}
\index{memory!distinction with respect to torsion}

The gravitational realization does not begin by identifying a discrete variable
with a geometric magnitude. It begins by distinguishing three levels: return
of observable transport, memory distinguishing histories with equal
projection, and holonomy of a connection. The bridge formulated here
receives the returns and cocycles already demonstrated in HMT; it does not reconstruct them
from the Holst action, from a mass, or from a metrological constant.

\section{Exact hierarchy of returns}
\label{sec:jerarquia-retornos-27-54-108}

Let \(F_{\rm obs}\) be the observable chronology defined in the HMT kernel, and let
\[
 \mathcal K_{27}:=F_{\rm obs}^{27}.
\]
Symbol \(\mathcal K_{27}\) designates exclusively the composition of
twenty-seven updates of the positional and oriented loop; it does not designate a
generic window, a trace of \(108\) steps, or the complete enriched TPK
state. Denote by \(p\) the projection onto the positions of the two
cursors and by \(\mathsf J_{\rm or}\) the simultaneous inversion of their
orientations. The primary proof is located in
\cref{prop:retorno-fobs-identidad}; its typed consequences are
\[
 p(\mathcal K_{27}x)=p(x),\qquad
 o(\mathcal K_{27}x)=\mathsf J_{\rm or}\,o(x),
 \qquad \mathsf J_{\rm or}^{\,2}=I,
\]
\[
 F_{\rm obs}^{54}:
 \text{restores positions and orientations},
 \qquad
 F_{\rm obs}^{108}=I:
 \text{also restores the observable phase}.
\]
Thus, position return, orientation return, and phase return are
distinct assertions. None of them by itself implies the return of
all fields of an enriched extension.

The mode quotient and orientation cocycles are established in
\cref{sec:cociclos-modo-orientacion,thm:descenso-necesario-fav}. In areal--volumetric
order,
\[
 N_{27}=9F_{\rm av},
 \qquad
 F_{\rm av}=
 \begin{pmatrix}0&1\\1&1\end{pmatrix}.
\]
If
\[
 \Omega_{\rm sw}(\gamma)=\sum_{a\subset\gamma}\omega_{\rm sw}(a),
 \qquad
 S(\gamma)=\sum_{a\subset\gamma}|\omega_{\rm sw}(a)|,
 \qquad
 G_{\rm or}(\gamma)=\sum_{a\subset\gamma}g(a),
\]
then a block of twenty-seven steps contains three nonadic turns and one
orientation reversal:
\[
 \Omega_{\rm sw}(\mathcal K_{27})=0,\qquad
 S(\mathcal K_{27})=18,\qquad
 G_{\rm or}(\mathcal K_{27})=1.
\]
In the observable return of \(108\) steps,
\[
 \Omega_{\rm sw}(C_{108})=0,\qquad
 S(C_{108})=72,\qquad
 G_{\rm or}(C_{108})=4.
\]
Thus, for the additive functional declared
\[
 \mathcal A_\ast(\gamma)
 :=G_{\rm or}(\gamma)+\lambda_\ast S(\gamma),
\]
are obtained, without introducing a dimensional scale,
\[
 \mathcal A_\ast(\mathcal K_{27})=1+18\lambda_\ast,
 \qquad
 \mathcal A_\ast(C_{108})=4+72\lambda_\ast.
\]
These values are record counts. Their subsequent thermal or gravitational
use does not alter their provenance.

\subsection{Directed system of coverings, refinement, and limit of the Cartan connection}
\label{sec:cobertura-acumulativa-dirigida}

Let \(\mathcal O_{108}\) be the observable orbit and, for each chronological edge
\(e\), let
\[
 c(e)=\bigl(g(e),s(e)\bigr)\in\mathbb N^2
\]
the pair formed by the indicators of orientation inversion and effective
sheet change. For a directed history
\(\gamma=e_1\cdots e_n\), define
\[
 c(\gamma)=\sum_{j=1}^{n}c(e_j).
\]
The concatenation law takes the form
\[
 c_{m+n}(x)
 =c_m(x)+c_n\!\left(F_{\rm obs}^{m}x\right).
\]
The previous counts express, uniformly over the orbit,
\[
 c_{27}(x)=(1,18),
 \qquad
 c_{108}(x)=(4,72).
\]

\begin{theorem}[Cumulative coverage of the observable chronology]
\label{thm:cobertura-acumulativa-dirigida}
On
\[
 \widetilde{\mathcal O}_{108}^{+}
 :=\mathcal O_{108}\times\mathbb N^2
\]
define
\[
 \widetilde F^{+}(x,m)
 :=\bigl(F_{\rm obs}x,m+c(x\to F_{\rm obs}x)\bigr).
\]
Then the transport of twenty-seven steps satisfies
\[
 \widetilde{\mathcal K}_{27}^{+}(x,m)
 =\bigl(\mathcal K_{27}x,m+(1,18)\bigr),
\]
and
\[
 \left(\widetilde{\mathcal K}_{27}^{+}\right)^4(x,m)
 =\bigl(x,m+(4,72)\bigr)\neq(x,m).
\]
Therefore, its observable projection has order four, whereas the
cumulative action of the chronological monoid has infinite semigroup order.
\end{theorem}

\begin{proof}
The cocyclic law and the constancy of \(c_{27}\) give
\[
 c_{108}(x)
 =\sum_{j=0}^{3}c_{27}(\mathcal K_{27}^{j}x)
 =4(1,18)=(4,72).
\]
The first coordinate returns through
\(\mathcal K_{27}^{4}=I\); the second does not return. After \(q>0\) iterations,
the increment is \(q(4,72)\), which is nonzero.
\end{proof}

\begin{remark}[Type of the cumulative extension]
The preceding construction extends directed histories through positive
counters. It is not an automorphic lift of groupoid
\(\mathcal G_{\rm mem}\), nor does it prove that the unbounded coordinate already belongs
to state \(\mathcal X_{\mathrm{TPK}}^{\mathrm{enr}}\). Its Grothendieck completion
\(\mathbb Z^2\) may be used for harmonic analysis, but it does not transform
the positive cost
\(\mathcal A_\ast=G_{\rm or}+\lambda_\ast S\) into a character of the isotropy
group: upon formally inverting a route, its indicators do not acquire
opposite signs.
\end{remark}

\paragraph{Domain of validity of the cumulative extension.}
The hierarchy \(27\to54\to108\), the incidences of \(F_{\rm av}\), and the
preceding cocycles are exact identities in the declared chronology and
readers. The assertion is restricted to the observable state preserved by
those readers; extending it to the full enriched state would require
proving the return of every forgotten field.

\section{Memory Groupoid}
\label{sec:grupoide-memoria-tpk}

Let \(\mathsf P_{\rm TPK}^{\rm enr}\) be the enriched path groupoid:
its objects are surviving TPK states, and its arrows are admissible routes
with declared source and target, composition by concatenation, and inversion
by route reversal. Let \(\sim_{\rm mem}\) be an equivalence relation that
identifies two arrows only when it preserves the memory fields required by
the reader under consideration and is congruent with identities, inversion, and
concatenation. We define
\[
 \mathcal G_{\rm mem}
 :=\mathsf P_{\rm TPK}^{\rm enr}/\!\sim_{\rm mem}.
\]
This definition does not transform memory into torsion. Memory belongs to
the discrete state and distinguishes histories; torsion belongs to a
geometric connection and measures \(D_\omega e\). Only a map between the two
domains can relate them.

A geometric realization requires, as additional data, a length
\(\ell_0>0\), a representation of the groupoid
\[
 \rho_{\rm C}:
 \mathcal G_{\rm mem}
 \longrightarrow
 \operatorname{Spin}(1,3)\ltimes\RR^{1,3},
\]
and a field map
\[
 \mathfrak R_{\rm grav}:
 \mathsf P_{\rm TPK}^{\rm enr}
 \dashrightarrow
 \mathsf{Cartan}_{1,3}.
\]
To define the second map, one must specify its action on objects and
routes, its compatibility with concatenation and inversion, and the regularity of the
resulting connection. None of these properties follows from the mere count
of returns.

\section{Joint connection and entanglement condition}
\label{sec:conexion-cartan-conjunta}

Under the geometric hypotheses of \cref{chap:gravedad} ---variedad
oriented of dimension four, co-tetrad \(e^I\), Lorentz connection
\(\omega^{IJ}\) and boundary conditions declared there---, the connection of
Cartan joint is written
\[
 \mathcal A_{\ell_0}
 =
 \begin{pmatrix}
  \omega&\ell_0^{-1}e\\
  0&0
 \end{pmatrix}.
\]
Its curvature is the algebraic identity.
\[
 \mathcal F_{\ell_0}
 =\dd\mathcal A_{\ell_0}
  +\mathcal A_{\ell_0}\wedge\mathcal A_{\ell_0}
 =
 \begin{pmatrix}
  F_\omega&\ell_0^{-1}T_{\rm tor}\\
  0&0
 \end{pmatrix},
\]
where
\[
 F_\omega^{IJ}
 =\dd\omega^{IJ}+\omega^I{}_K\wedge\omega^{KJ},
 \qquad
 T_{\rm tor}^{I}=D_\omega e^I.
\]
Curvature and torsion appear as distinct components of the same
transport defect. The matrix equality follows algebraically from the
explicit data \(e,\omega,\ell_0\), but does not construct these inputs from TPK.

The strong content of the bridge is the commutativity of holonomies:
\[
 \boxed{
 \Hol_{\mathcal A_{\ell_0}}
 \bigl(\mathfrak R_{\rm grav}(\gamma)\bigr)
 =
 \rho_{\rm C}\!\left(
 \Hol_{\rm TPK}(\gamma)
 \right)}
 \qquad
 (\gamma\in\mathsf P_{\rm TPK}^{\rm enr}).
\]
The equality must be compatible with identities, inversion, concatenation,
subdivision, change of base point, and conjugation. Only then is discrete
memory represented in a connection; it is not merely renamed as
torsion or curvature.

\begin{criterion}[Existence conditions for the bridge TPK--Cartan]
\label{crit:puente-tpk-cartan}
There is not realization with the declared commutativity and compatibility
if any of the following facts occurs:
\begin{enumerate}
 \item two equivalent representatives of the same route produce
 nonconjugate geometric holonomies;
 \item the image does not preserve identities, inverses, or concatenation;
 \item an admissible subdivision alters the holonomy;
 \item the equality of holonomies fails for \(\mathcal K_{27}\) or for their
 concatenations of depths \(54\) and \(108\);
 \item the map identifies memory and torsion without exhibiting a morphism that
 preserves their types;
 \item some descending choice ---the Holst action, \(G\), a mass or
 an SI chart--- is used to redefine \(A\), \(C^\ast\) or the original
 TPK return.
\end{enumerate}
\end{criterion}

\paragraph{Antecedence of the functional \(\Gamma_{6\to8}\) with respect to the Palatini--Cartan--Holst realization.}
\cref{mdg:chap:barbero} fixes, once \(A,C^\ast,q_\pm\) have been constructed, the
internal functional \(\Gamma_{6\to8}\). \cref{chap:gravedad} then studies
the Palatini--Cartan--Holst action and its variational equations. The present
bridge does not introduce a reverse route between the two stages.
